\documentclass[12pt]{article}
\usepackage[utf8]{inputenc}
\usepackage[T1]{fontenc}
\usepackage{hyperref}
\title{Mohegan-Pequot Language Resources for K-12 Curriculum Development}
\author{Compiled by Claude Sonnet 5.5 (claude-sonnet-5-5) for Andr\'as Kornai}
\date{}
\begin{document}
\maketitle

\section{Introduction}
Mohegan-Pequot (ISO 639-3 xpq) is an Eastern Algonquian language of southern New England and eastern Long Island. Its last fluent speaker, Fidelia Fielding, died in 1908. Her diaries, written in Mohegan, were preserved and later repatriated by the tribe. They are central to the revival. Stephanie Fielding, a Mohegan linguist, compiled \emph{A Modern Mohegan Dictionary}. My searches found only general and partly dated sources on current programmes, so the status of school instruction should be checked with the Mohegan Tribe and the Mashantucket Pequot Tribal Nation directly.

\section{Primary Educational Resources}
I found no primers, textbooks or K-12 curricula in the results.

\section{Linguistic and Reference Materials}
\begin{itemize}
\item Wikipedia, ``Mohegan-Pequot language''. \url{https://en.wikipedia.org/wiki/Mohegan-Pequot_language}
\item Wikipedia, ``Stephanie Fielding'': the Mohegan linguist who compiled \emph{A Modern Mohegan Dictionary}. \url{https://en.wikipedia.org/wiki/Stephanie_Fielding}
\item FreeLang Mohegan dictionary page. \url{https://www.freelang.net/dictionary/mohegan.php}
\item Wikipedia, ``Mohegan Tribe''. \url{https://en.wikipedia.org/wiki/Mohegan_Tribe}
\end{itemize}

\section{Religious and Community Resources}
\begin{itemize}
\item Indianz.com 2004 news report on language revitalization. It said the Mashantucket Pequot Tribal Nation had adopted a strategic plan and funding, while the Mohegan Tribe was less far along. The report is dated. \url{https://www.indianz.com/News/2004/000989.asp}
\item Brothertown Indians page, surfaced in the search; I did not verify its content. \url{https://brothertownindians.org/?p=266}
\end{itemize}

\section{Digital and Audio Resources}
The FreeLang page above is the only online word list found. No audio was found in the results.

\section{Curriculum Development Context}
The Mohegan revival rests on Fidelia Fielding's diaries and Stephanie Fielding's dictionary, so school materials would need to be built with tribal language staff. Sources disagree on when the revival began, and one report notes that the Wôpanâak Language Reclamation Project has influenced Eastern Algonquian revival methods. Current tribal language programmes should be confirmed with the tribes before planning.
\end{document}